\documentclass[10pt,letterpaper]{article}
\usepackage{cornell-notes}

\title{Annex C.07.01: 07-01-general}
\author{}
\date{}
\setDocTitle{Annex C.07.01: 07-01-general}
\setDocSubtitle{ISO/IEC/IEEE 42010:2022}
\setDocOwner{ISO 42010 Cornell Notes}
\setDocDate{\today}
\setCornellCollection{ISO/IEC/IEEE 42010:2022 Cornell Notes}
\setCornellUnitType{Annex}
\setCornellUnitNumber{C.07.01}
\setCornellUnitTitle{07-01-general}
\hypersetup{pdftitle={Annex C.07.01: 07-01-general},pdfsubject={Cornell notes},pdfauthor={}}

\begin{document}
\maketitle
\begin{CornellOverviewBox}{Overview}
{\Large\bfseries\textcolor{CNavy}{Annex C.7.1: General}}\\[3pt]
{\small\textbf{Classification:} Informative annex}\\
{\small\textbf{Learning objective:} Use the informative guidance on general to reinforce the normative and conceptual material without treating the guidance itself as mandatory.}
\end{CornellOverviewBox}
\vspace{3pt}

\begin{CornellNotesTable}
\CornellNoteRow{Purpose / key concept}{The RM-ODP framework defines five viewpoints for specifying ODP systems and a set of correspondences between them.}
\CornellNoteRow{Core details}{\begin{itemize}[leftmargin=*,nosep,topsep=1pt]
\item For each viewpoint, there is an associated viewpoint language which defines “the concepts and rules for specifying ODP systems from the corresponding viewpoint”.
\item An AD conforming to This section and using ISO/IEC 10746-3 would include the viewpoints defined by ISO/IEC 10746-3 and views to implement these viewpoints.
\item A conforming AD does not need to be limited to the five predefined viewpoints of ISO/IEC 10746-3; the AD can include additional viewpoints and views, as needed.
\item Elements of that specification specific to ADs (such as stakeholders) are omitted here since they are particular to individual systems.
\end{itemize}}
\CornellNoteRow{Interpretive note}{\begin{itemize}[leftmargin=*,nosep,topsep=1pt]
\item ISO/IEC 19793 defines a UML profile for the specification of open distributed processing systems using these viewpoints.
\end{itemize}}
\CornellNoteRow{Active recall}{\begin{itemize}[leftmargin=*,nosep,topsep=1pt]
\item What is the central role of General?
\item How does General connect to adjacent architecture-description concepts?
\end{itemize}}
\end{CornellNotesTable}


\begin{CornellSummaryBox}{Summary}
The RM-ODP framework defines five viewpoints for specifying ODP systems and a set of correspondences between them. This material is informative guidance and does not itself create a conformance requirement.
\end{CornellSummaryBox}

\end{document}
